\documentclass[11pt,a4paper]{article}
\usepackage[margin=2.1cm]{geometry}
\usepackage[T1]{fontenc}
\usepackage{lmodern}
\usepackage{microtype}
\usepackage{booktabs}
\usepackage{tabularx}
\usepackage{array}
\usepackage{enumitem}
\usepackage[colorlinks=true,urlcolor=blue]{hyperref}

\hyphenpenalty=10000
\exhyphenpenalty=10000
\emergencystretch=3em
\linespread{1.02}
\setlist{itemsep=3pt,topsep=4pt,parsep=0pt}
\setlength{\parindent}{0pt}
\setlength{\parskip}{4pt}
\renewcommand{\arraystretch}{1.15}

\title{RAG Assignment Report: Baseline Retrieval vs Cross Encoder Reranking}
\author{MD TAUFIQUE HUSSAIN \quad Roll No. 2025202007}
\date{GitHub: \url{https://github.com/Taufiquesir/rag-baseline-vs-reranking}}

\begin{document}
\maketitle
\vspace{-1.2em}

\section{Introduction and setup}

This report focuses on two different approaches to retrieve 3 chunks from the technical PDF to pass to the LLM in a small scale RAG architecture. For the Pipeline 1 (Baseline), a Bi Encoder finds out the top 10 chunks according to the cosine similarity score, then the top 3 among them are passed to the LLM and Answer 1 is generated. In the Pipeline 2 (Reranking), 10 chunks are ranked once again using a Cross Encoder based on their scores, then the top 3 chunks are passed to the LLM to get Answer 2. The chunks, the 10 candidate chunks, the LLM and the prompt are the same in both cases.
\begin{itemize}
\item \textbf{Document.} NIST SP 800-63B-4, Digital Identity Guidelines: Authentication and Authenticator Management (July 2025, 129 pages), the US standard for passwords, multifactor login, one time codes, passkeys and sessions. It was chosen because it should be hard for plain vector search: many rules use SHALL NOT, many rules have an exception (``unless'', ``however''), the paragraphs for assurance levels AAL1, AAL2 and AAL3 are nearly identical and differ mainly in a number, and the standard says ``PSTN'' where a user would say ``SMS''.
\item \textbf{Preprocessing.} The notebook deletes the initial 11 pages (title, legal statements, table of contents), as well as the header at the top of each page, and after that, splits the text into chunks of 120 words with an overlap of 30 words. This gives 413 chunks.
\item \textbf{Models.} Bi Encoder \texttt{all-MiniLM-L6-v2} with exact FAISS cosine search, Cross Encoder \texttt{ms-marco-MiniLM-L-6-v2}, and the LLM Qwen2.5 1.5B Instruct on a laptop CPU with greedy decoding. The prompt tells the LLM to answer only from the given passages or reply ``Not found''.
\item \textbf{Checking retrieval and answers.} 
The gold chunk that has the precise sentence answering the test questions is identified for each test question and its rank is noted for both pipelines. Ranks up to 3 can only make it to the LLM. The answers are compared manually to the standard sentence, and if the rule is stated correctly in the answer, then it is marked as correct.
\end{itemize}

\section{Architecture comparison}

The Bi Encoder processes the question and each chunk independently. Then, it converts each to a vector and compares two vectors using the cosine similarity measure. Since a chunk vector is independent from the question, all 413 chunk vectors were precomputed (13.5 seconds) and it takes only 18 ms to encode a question and compare all vectors. The downside is that the entire chunk is compressed into a vector without taking the question into account, which may result in loss of minor details like a number, a level name (e.g., AAL1), or the word ``not''. That was the case for query 6 where the level AAL1 is required. The paragraphs of the AAL2 level (chunk 26, cosine 0.752) and AAL3 level (chunk 31, 0.747) scored higher than the AAL1 paragraph (chunk 20, 0.731).

Cross Encoder sees the question and the chunk in combination as the input and self attention makes each word of the question to attend to every word of the chunk. This is a possible explanation why Cross Encoder could consider the AAL1 paragraph as first in ranking, but this experiment does not identify the reason. The disadvantage of this approach is that no preparation is allowed; it has to work for each question and chunk combination, and for each combination it takes around 37 ms on the laptop in use. That is why these two models are used in conjunction. The quick Bi Encoder selects 10 candidates from the entire document, and Cross Encoder ranks these 10.

\section{Qualitative and rank shift analysis}

Reranking put the gold chunk in the top 3 in all 8 queries, so that retrieval became successful for all eight cases. However, not all of the final answers were successful for all eight. For the context retrieved with reranking, the LLM provided the correct answer for 6 of 8 queries (but not for queries 2 and 4). In the case of baseline context, none of the 8 answers was completely correct, and only the query 3 had the correct conclusion, but with incorrect reasoning. The table shows the ranks of the gold chunk and LLM answers. Queries 9 to 11 serve as counter examples.

{\small
\newcolumntype{R}[1]{>{\raggedright\arraybackslash}p{#1}}
\begin{tabularx}{\textwidth}{@{}r R{1.9cm} >{\raggedright\arraybackslash}X c c R{2.6cm} R{3.1cm}@{}}
\toprule
\# & Type & Question & Base & Rerank & Answer 1 & Answer 2 \\
\midrule
1 & Exception & When must a user be forced to change a password? & \#5 & \#1 & unclear, only points to a section & correct, only after compromise \\
2 & Exception & Is encryption required for one time codes sent by SMS? & \#10 & \#1 & not found & wrong, missed the PSTN exception \\
3 & Exception & Can synced passkeys be used at AAL3? & \#7 & \#3 & right answer, wrong reason & correct, keys must not be exportable \\
4 & Negation & Must users change passwords every 90 days? & \#5 & \#1 & not found & not found \\
5 & Negation & Does a new OTP reset the failed attempt counter? & \#5 & \#1 & not found & correct, SHALL NOT reset \\
6 & Keyword trap & What is the reauthentication timeout at AAL1? & \#6 & \#1 & not found & correct, 30 days \\
7 & Keyword trap & What is the inactivity timeout at AAL3? & \#4 & \#1 & wrong, 1 hour & correct, 15 minutes \\
8 & Synonym & Is the use of SMS for login restricted? & \#4 & \#1 & confused & correct, PSTN is restricted \\
9 & Counter & Can a fingerprint alone be used to log in? & \#3 & \#7 & no, wrong reason & no, wrong reason \\
10 & Counter & What is the inactivity timeout at AAL2? & \#2 & \#3 & correct & correct \\
11 & Counter & Can email be used to send one time login codes? & none & none & not found & hallucinated yes \\
\bottomrule
\end{tabularx}
}

\textbf{Clearest case, query 7.} The baseline top 3 were a general session section (chunk 218) and two chunks from the AAL2 text (chunks 25 and 26), which give the inactivity timeout as no more than 1 hour. The AAL3 paragraph (chunk 31) was at rank 4, just outside the top 3, so Answer 1 gave the AAL2 value and misstated the AAL3 rule. The Cross Encoder moved chunk 31 to rank 1, with a score of 6.79 against 6.43 for the best AAL2 chunk, and Answer 2 said 15 minutes. Query 6 shows the same problem with a smaller margin. The top six baseline candidates (cosine 0.731 to 0.771) all came from the sections on assurance levels and reauthentication, and the Cross Encoder put the AAL1 chunk first with a score of 7.18 against 7.12 for the AAL3 chunk.

\textbf{Other retrieval fixes.} In query 1, the top 3 baseline chunks concerned password creation and modification (a blocklist check and two usability chunks). The rule was at rank 5, and Answer 1 only pointed to Sec. 3.1.1 without stating the rule. In query 3, two out of three baseline chunks were taken from Appendix B and concerned synced authenticators, while the last one was related to exportable keys at AAL2 or lower. The rule banning synced passkeys at AAL3 was at rank 7. In query 8, the top baseline chunk was a usability chunk with reference to SMS solely as an example of a suitability problem; the other two were also usability chunks. The chunk defining PSTN as the only restricted authenticator was at rank 4.

\textbf{Better retrieval was not always enough.} In Query 2, even though the gold chunk shifted from position 10 to position 1, the LLM stated that encryption was necessary and removed the ``unless sent via the PSTN'' exception, even though the full sentence appeared in Passage 1. In Query 4, the rule was ranked \#1 and its sentence was at the start of Passage 1, yet the LLM returned ``Not found''. One possible explanation here is that the question uses ``every 90 days,'' whereas the rule states ``periodically,'' but this has not been confirmed experimentally.

\textbf{Counter examples.} The Cross Encoder moved chunk 121, which stated ``Biometrics SHALL only be used as part of multifactor authentication'', from position 3 to position 7 in query 9. Its top three choices were a usability table, an introductory paragraph for the biometrics section and a chunk regarding the difficulty of collecting fingerprints. Its highest score was minus 7.62, with all ten of its scores being equal to or lower than this figure, indicating that it was not very confident in its candidates. Query 11 involved a gold chunk (chunk 69, the same chunk that answers query 2), which was not in the top 10 chunks, so reranking could not help. The Cross Encoder's top choice for this query was a chunk discussing the issuing of recovery codes at a recovery address (postal, email, text or voice). Answer 2 said yes, so a safe ``Not found'' became a wrong answer.

\textbf{Wider test.} Retrieval alone was also run on 48 questions. These are all the questions that were tried while looking for failure cases, so they are not a random sample and the scores should not be read as typical performance.

\begin{center}
\small
\begin{tabular}{@{}lccc@{}}
\toprule
 & Hit@1 & Hit@3 & MRR@10 \\
\midrule
Pipeline 1, Bi Encoder & 0.250 & 0.583 & 0.437 \\
Pipeline 2, plus Cross Encoder & 0.500 & 0.750 & 0.617 \\
\bottomrule
\end{tabular}
\end{center}

The gold rank improved in 19 questions, stayed the same in 15 and became worse in 5. In the remaining 9 questions (19 percent) the gold chunk was never in the top 10, so Hit@10 is 0.812 for both pipelines. Five of these nine gold sentences are SHALL NOT rules (composition rules, password hints, voice biometrics, email for out of band codes, and manual entry codes not being phishing resistant). Reranking also helped with 100 word chunks (Hit@3 0.542 to 0.688) and 150 word chunks (0.583 to 0.708), so the result does not depend on the chunk size.

\section{Latency profiling}

Each step was timed with \texttt{time.perf\_counter()}, using 20 runs for each of the 11 queries after a warm up, on a laptop CPU (Ryzen 5 5500U, no GPU). The LLM time was measured once per answer.

\begin{center}
\small
\begin{tabular}{@{}lrrr@{}}
\toprule
Step & Mean & Median & 95th percentile \\
\midrule
Pipeline 1: Bi Encoder retrieval (top 10) & 17.8 ms & 17.5 ms & 23.2 ms \\
Pipeline 2: extra Cross Encoder step (10 pairs) & 365.7 ms & 350.4 ms & 489.2 ms \\
Pipeline 2: total retrieval & 383.5 ms & 368.1 ms & 509.4 ms \\
LLM generation, Answer 1 & 51.0 s & 50.5 s & 60.2 s \\
LLM generation, Answer 2 & 51.7 s & 51.9 s & 56.4 s \\
End to end, Pipeline 1 & 51.02 s & 50.51 s & 60.19 s \\
End to end, Pipeline 2 & 52.08 s & 52.25 s & 56.82 s \\
\bottomrule
\end{tabular}
\end{center}

Reranking raises retrieval time from 17.8 ms to 383.5 ms, about 21.5 times slower, because the Cross Encoder runs the transformer on 10 pairs while the Bi Encoder only encodes the question. The extra 366 ms is 0.7 percent of the baseline LLM time of about 51 seconds. The measured end to end time rose from 51.02 s to 52.08 s, about 2.1 percent. Reranking explains 0.37 s of this 1.06 s rise, and the other 0.70 s is the difference between the two LLM rows. LLM times varied a lot between questions (for Answer 1 the 95th percentile is 60.2 s against a median of 50.5 s) and each answer was timed once, so a 0.7 s difference is inside this spread and the data cannot assign it to reranking.


\section{Scaling question}

Cross Encoder cannot be applied across the entire document since it requires the question along with the chunk; therefore, nothing can be computed until the question comes. It uses the transformer N times on N chunks for each question asked, and thus the cost increases linearly with N. This was evaluated for Query 1: scoring all 413 chunks took 16.9 seconds (around 41 milliseconds per pair, marginally higher than 37 milliseconds per pair mentioned above for a different set of computations), which is 44 times slower compared to the two stage model that took 0.38 seconds. In this case, the gold chunk scored the highest for both, and thus the extra time did not produce any advantage, although one query is not enough to say this holds in general.

\begin{center}
\small
\begin{tabular}{@{}rrr@{}}
\toprule
Number of chunks & Cross Encoder on every chunk & Two stage pipeline \\
\midrule
413 (measured) & 16.9 s & 0.38 s \\
10,000 (estimated) & 6.8 minutes & 0.38 s \\
1,000,000 (estimated) & 11.3 hours & 0.38 s \\
\bottomrule
\end{tabular}
\end{center}

The Bi Encoder computes every chunk vector once, offline. At question time it encodes only the question and searches the stored vectors, and an approximate index such as HNSW keeps this in milliseconds even for millions of chunks. The Cross Encoder then sees only 10 candidates, so its cost stays the same however large the collection grows.

\section{Conclusion and limitations}

Reranking helped in selecting the right chunks for the LLM in each of the failing queries and increased the Hit@3 on the
48 questions from 0.583 to 0.750, at a cost of about 0.37 seconds per question. The improvement showed up where the
baseline had to choose between many similar candidates (AAL level) or when the answer was a negation or exception. The limitations are clear. The gold rank got worse in 5 of 48 questions, it cannot compensate for the 9 questions where stage 1 did
not find the gold chunk, and the LLM failed in queries 2 and 4 even with the right text in front of it. Next I would try
retrieval of more candidates (top 30 to 50), BM25 keyword search (``email'', ``PSTN'') and a bigger LLM.
The limitations of this study are that the 48 questions and gold sentences were written specifically for this assignment and not taken from a benchmark,
the answers were evaluated manually, the LLM is small, the LLM time for each example was measured only once, the
whole document scan took place for one query and all timings are from one laptop.

\end{document}
